% Nouveau EGA V (cote25): archive cote25, site120 (PDF page121).
% English translation of the checked French annotated master.
% AI-translated by GPT-6 Astra (OpenAI, ChatGPT Pro mode).
\scanpage{25/120}
\mspage{\unsure{23|27}}

Z$_s$ has codimension $\hand{\geq} 2$ in X$_s$ ; \var{\xout{or again (as in the cases consid}}{}{typewriter strike-out}{ega25_25_120-v0001}
$\hand{\unsure{\phi|\emptyset}}$ is then the set of \var{\xout{families}}{}{typewriter strike-out}{ega25_25_120-v0002} \var{\inserted[typed]{lin syst}}{lin syst}{insertion}{ega25_25_120-v0003} of divisors param.by P which \unsure{are}
\uline{without fixed component}, \var{\xout{whereas $\hand{\unsure{\psi|\emptyset}}$ is the set of}}{}{typewriter strike-out}{ega25_25_120-v0005} On the other hand, if
S=Spec(k), and X is normal, then \var{\inserted[typed]{for}}{for}{insertion}{ega25_25_120-v0006} \uline{every} rational map
f:X$\hand{\to}$P, the domain of definition U(f) satisfies codim ((X-U(f),X)$\hand{\geq}$2
(II 7), so in this case $\hand{\unsure{\psi|\emptyset}}$ consists of the set of \uline{all}
rational maps from X into P.

The proof of a) has already been given. To prove b), note
that formation of $i_{\hand{*}}(\underline{L}_U)$ commutes with every flat extension of the
\var{\xout{for the\unsure{x} one must}}{}{typewriter strike-out}{ega25_25_120-v0007}
base S' \quad S (ref)\var{\inserted{at least when $U\to X$ is quasi-compact, a case to which one reduces without \unsure{difficulty}}}{at least when $U\to X$ is quasi-compact, a case to which one reduces without \unsure{difficulty}}{insertion}{ega25_25_120-v0008}, hence \var{\inserted[typed]{the cond. in question is inv. under}}{the cond. in question is inv. under}{insertion}{ega25_25_120-v0009} \var{\xout{it suffices to find a morphism}}{}{typewriter strike-out}{ega25_25_120-v0010}
\var{\xout{S' $\quad$ S such that the conclusion is true after}}{}{typewriter strike-out}{ega25_25_120-v0011} \var{\inserted[typed]{faith flat qu cpc}}{faith flat qu cpc}{insertion}{ega25_25_120-v0012} base change. One
will take S' = $\hand{\check{}}\mathrm{P}$ , and note that the hypothesis that $i'_{\hand{*}}(\underline{L}'_U)$ is invertible does not
change on replacing $\underline{L}'_U$ by $\underline{L}'_U\otimes_{S'}M'$, where M' is an invertible Module
on S' (since $i'_{\hand{*}}(\underline{L}'_U\otimes_{S'}M')\hand{\simeq} i'_{\hand{*}}(\underline{L}'_U)\otimes_{S'}M'$ ). One will \unsure{take} M'= $\underline{O}_{\overset{\hand{\vee}}{P}}(1)$, so
that the condition in question also means that $i'_{\hand{*}}(\underline{N}')$ is an invertible
Module, where $\underline{N}'= (f_U\times_S id_P)^{\hand{*}}(\underline{O}_{P\times\overset{\hand{\vee}}{P}}(1,1))$ . Now $\underline{O}$ (1,1) is precisely
the invertible Module defined by the canonical divisor $\underline{H}$ on $P\times_S\overset{\hand{\vee}}{P}$ , so
that $\underline{N}$ is just the invertible Module defined by \qquad D' =
$(f_U\times id_P)$ $(\underline{H})$. \var{\deleted{\var{\xout{\illegible{1}\unsure{D N P}\illegible{1}\unsure{Now}}}{}{typewriter strike-out}{ega25_25_120-v0014}\unsure{S}}}{}{deletion}{ega25_25_120-v0013} When f comes from D$\hand{\in}\hand{\unsure{\phi|\emptyset}}$ , \var{\xout{the latter}}{}{typewriter strike-out}{ega25_25_120-v0015}
D' is just D$\hand{|}$U$\times_S\overset{\hand{\vee}}{P}$ , \var{\xout{and it follows from the lemma}}{}{typewriter strike-out}{ega25_25_120-v0016} hence $\underline{N}' = \underline{N}\hand{|}U\times_S\overset{\hand{\vee}}{P}$ ,
where $\underline{N}$ is the invertible Module on $X\times_S\hand{\check{}}P$ defined by D, and it follows
from lemma \var{\inserted{3.2}}{3.2}{insertion}{ega25_25_120-v0017} above \var{\xout{that one will have}}{}{typewriter strike-out}{ega25_25_120-v0018} applied to $P\times_S\hand{\check{}}P\hand{\to}\hand{\check{}}P$ that one will have
$i'_{\hand{*}}(\underline{N}')\hand{\simeq}\underline{N}$ , so $i'_{\hand{*}}(\underline{N}')$ is invertible. Conversely, when this condition
holds, let us prove that f comes from a D$\hand{\in}\hand{\unsure{\phi|\emptyset}}$ , or, manifestly
equivalently, that the divisor D' extends to a \var{\deleted{Cartier}}{}{deletion}{ega25_25_120-v0022}
divisor relative \var{\xout{of}}{}{typewriter strike-out}{ega25_25_120-v0023} to $\hand{\check{}}P$ on X$\times_S$P . It \var{\deleted{suffices}}{}{deletion}{ega25_25_120-v0024} amounts to the same
to say that it extends to a \var{\deleted{Cartier}}{}{deletion}{ega25_25_120-v0025} divisor D on X$\times_S$P , for \hand{D}.

\var{\margin{\var{\deleted{\unsure{Complicated remark}}}{}{deletion}{ega25_25_120-v0027} \unsure{One} must \unsure{express|explain} \illegible{1} $(f_U\times_S\mathrm{id}_P)^{-1}(H)$ \unsure{which} is defined as a \unsure{divisor} rel. to P \unsure{thanks} to the condition \illegible{3} \unsure{$f_U$} \unsure{extends to} \unsure{a} relative divisor /P \unsure{on} $X\times_S P$ \unsure{/P} \var{\deleted{\illegible{2}}}{}{deletion}{ega25_25_120-v0028} \unsure{*}}}{\placedmargin{\var{\deleted{\unsure{Complicated remark}}}{}{deletion}{ega25_25_120-v0027} \unsure{One} must \unsure{express|explain} \illegible{1} $(f_U\times_S\mathrm{id}_P)^{-1}(H)$ \unsure{which} is defined as a \unsure{divisor} rel. to P \unsure{thanks} to the condition \illegible{3} \unsure{$f_U$} \unsure{extends to} \unsure{a} relative divisor /P \unsure{on} $X\times_S P$ \unsure{/P} \var{\deleted{\illegible{2}}}{}{deletion}{ega25_25_120-v0028} \unsure{*}}}{marginal note placed}{ega25_25_120-v0026}
\tnote{25-120-1}{The inserted doubtful \unsure{ie} changes French ``sont'' to ``soient'', both rendered ``are'' here. The doubt in ``take'' concerns French ``prèndra|prendra''. The first canceled fragment ends at ``envisa'', retained as ``consid''. ``Param.by'' renders ``par.par'', parametrized by; ``lin syst'', ``cond.'', ``inv.'', and ``rel.'' mean linear system, condition, invariant, and relative. ``Faith flat qu cpc'' retains the abbreviated expression ``fid plat qu cpc'', faithfully flat and quasi-compact. The canceled marginal ``Rq'' means remark. The source's incomplete final ``D.'' and the differing projective-space symbols remain.}

\mspage{120}
